% Figure: blackbody spectral emissive power (Planck) vs wavelength for a
% few temperatures, illustrating the Wien shift of the peak to shorter
% wavelength as T rises and the strong growth of total radiated power.
% Curves are normalised per curve to peak near unity for display on one
% set of axes; the normalisation constants are the Planck peak values.
\begin{figure}[htbp]
\centering
\begin{tikzpicture}
\begin{axis}[
  width=12cm, height=7.5cm,
  xlabel={wavelength $\lambda$ ($\mu$m)},
  ylabel={spectral emissive power (normalised)},
  xmin=0, xmax=12, ymin=0, ymax=1.15,
  axis lines=left,
  legend pos=north east,
  legend cell align=left,
  tick label style={font=\scriptsize},
  label style={font=\small},
  legend style={font=\scriptsize},
  clip=false,
]
% Planck shape E_lambda ~ 1/(lambda^5 (exp(c2/(lambda T)) - 1)), c2=14388
% um*K, each curve normalised to its own peak. Coordinates precomputed to
% keep pgfmath out of its fixed-point overflow range (lambda^5 overflows).
\addplot[engred, very thick] coordinates {
(0.40,0.0000) (0.80,0.0000) (1.00,0.0014) (1.20,0.0110) (1.40,0.0435)
(1.60,0.1111) (1.80,0.2150) (2.00,0.3449) (2.20,0.4851) (2.40,0.6207)
(2.60,0.7406) (2.80,0.8386) (3.00,0.9122) (3.20,0.9619) (3.40,0.9901)
(3.60,0.9999) (3.80,0.9946) (4.00,0.9774) (4.40,0.9185) (4.80,0.8416)
(5.20,0.7585) (5.60,0.6766) (6.00,0.5996) (6.60,0.4971) (7.20,0.4111)
(8.00,0.3198) (9.00,0.2357) (10.00,0.1761) (11.00,0.1335) (12.00,0.1026)
};
\addlegendentry{$T=800$ K (peak $\approx 3.6\,\mu$m)}
\addplot[engamber, very thick] coordinates {
(0.40,0.0000) (0.60,0.0003) (0.80,0.0110) (1.00,0.0725) (1.20,0.2150)
(1.40,0.4147) (1.60,0.6207) (1.80,0.7926) (2.00,0.9122) (2.20,0.9785)
(2.40,0.9999) (2.60,0.9873) (2.80,0.9512) (3.00,0.9004) (3.40,0.7794)
(3.80,0.6568) (4.20,0.5463) (4.60,0.4521) (5.00,0.3740) (5.60,0.2826)
(6.20,0.2156) (7.00,0.1530) (8.00,0.1026) (9.00,0.0710) (10.00,0.0504)
(11.00,0.0368) (12.00,0.0274)
};
\addlegendentry{$T=1200$ K (peak $\approx 2.4\,\mu$m)}
\addplot[englime!70!black, very thick] coordinates {
(0.40,0.0003) (0.60,0.0324) (0.80,0.2150) (1.00,0.5199) (1.20,0.7926)
(1.40,0.9516) (1.60,0.9999) (1.80,0.9716) (2.00,0.9004) (2.20,0.8107)
(2.40,0.7171) (2.60,0.6277) (2.80,0.5463) (3.00,0.4741) (3.40,0.3567)
(3.80,0.2700) (4.20,0.2064) (4.60,0.1595) (5.00,0.1248) (5.60,0.0882)
(6.20,0.0639) (7.00,0.0429) (8.00,0.0274) (9.00,0.0182) (10.00,0.0126)
(11.00,0.0089) (12.00,0.0065)
};
\addlegendentry{$T=1800$ K (peak $\approx 1.6\,\mu$m)}
\node[font=\scriptsize, engblue, anchor=west] at (axis cs:4.6,0.92)
  {Wien: $\lambda_{\max}T=2898\,\mu\mathrm{m\cdot K}$};
\draw[engblue, dashed, ->] (axis cs:4.55,0.92) -- (axis cs:1.7,0.95);
\end{axis}
\end{tikzpicture}
\caption[Blackbody spectral emissive power]{Blackbody spectral emissive
power against wavelength for three temperatures, each curve normalised to
its own peak so all three fit on one set of axes. The peak wavelength
moves to the left as temperature rises, following the Wien displacement
law $\lambda_{\max}T = 2898\,\mu\mathrm{m\cdot K}$: a body at $800$ K
radiates mostly in the infrared, while at $1800$ K an appreciable part of
the emission has crossed into the near infrared and the body begins to
glow visibly. The total area under each unnormalised curve grows as
$T^4$, the Stefan-Boltzmann result, so the hotter body radiates far more
power than the normalised plot suggests.}
\label{fig:vol04ch05:blackbody-spectrum}
\end{figure}
